\section{Agentic Systems}
\label{sec:agentic-systems}

Các chương trước tập trung vào jailbreak attacks, defenses và evaluation
đối với LLM sinh văn bản. Tuy nhiên, khi LLM được tích hợp vào hệ thống có
bộ nhớ, khả năng lập kế hoạch, truy xuất dữ liệu và gọi công cụ, bài toán
an toàn không còn chỉ nằm ở phản hồi cuối cùng. Một agent có thể đọc dữ
liệu từ nguồn bên ngoài, ra quyết định nhiều bước, gọi API, ghi bộ nhớ hoặc
tương tác với môi trường. Vì vậy, rủi ro jailbreak và prompt injection cần
được phân tích ở cấp toàn bộ trajectory của hệ thống.

Chương này trình bày tổng quan về LLM agents, bề mặt tấn công của agent,
prompt injection trực tiếp và gián tiếp, tool-use attacks, memory poisoning,
jailbreaking đối với robot được điều khiển bởi LLM, các phương pháp phòng
thủ đặc thù cho agent và benchmark cho agent security.

\subsection{Tổng quan LLM agents}
\label{subsec:llm-agents-overview}

Một chatbot thông thường nhận prompt từ người dùng và sinh phản hồi văn
bản. Ngược lại, một LLM agent sử dụng mô hình ngôn ngữ như bộ điều khiển
cho một vòng lặp hành động. Agent không chỉ trả lời câu hỏi mà còn có thể
lập kế hoạch, gọi công cụ, quan sát kết quả và điều chỉnh bước tiếp theo.

Một agent đơn giản có thể được mô tả bởi vòng lặp sau:

\begin{equation}
    o_t \rightarrow p_t \rightarrow a_t \rightarrow o_{t+1},
    \label{eq:agent-loop}
\end{equation}

trong đó $o_t$ là observation tại thời điểm $t$, $p_t$ là kế hoạch hoặc
suy luận trung gian, $a_t$ là hành động được chọn và $o_{t+1}$ là quan sát
mới sau khi hành động được thực thi.

Trong thực tế, một agent thường gồm các thành phần chính:

\begin{itemize}[leftmargin=1.2cm]
    \item \textbf{LLM controller:}
    mô hình ngôn ngữ chịu trách nhiệm diễn giải nhiệm vụ, lập kế hoạch và
    chọn hành động;

    \item \textbf{Memory:}
    bộ nhớ ngắn hạn hoặc dài hạn dùng để lưu thông tin từ các lượt tương
    tác trước;

    \item \textbf{Planner:}
    thành phần tạo kế hoạch nhiều bước hoặc phân rã nhiệm vụ;

    \item \textbf{Tools:}
    các công cụ bên ngoài như search, database, email, file system, code
    interpreter hoặc API ứng dụng;

    \item \textbf{Environment:}
    môi trường mà agent quan sát và tác động, chẳng hạn website, tài liệu,
    hệ thống phần mềm hoặc robot;

    \item \textbf{Guard/monitor:}
    các lớp kiểm soát an toàn, logging và xác nhận hành động.
\end{itemize}

Chính việc bổ sung memory, tools và environment làm tăng đáng kể bề mặt tấn
công. Trong chatbot, attacker chủ yếu tác động qua prompt trực tiếp. Trong
agentic system, attacker còn có thể tác động qua dữ liệu được retrieval,
tool output, memory entry, trạng thái website hoặc môi trường vật lý.

\subsection{Bề mặt tấn công của agent}
\label{subsec:agent-attack-surface}

Bề mặt tấn công của agent rộng hơn chatbot vì agent xử lý nhiều nguồn dữ
liệu và thực hiện nhiều loại hành động hơn. Thay vì chỉ đánh giá một phản
hồi văn bản, cần xét toàn bộ pipeline từ prompt đầu vào đến các bước lập
kế hoạch, truy xuất dữ liệu, gọi công cụ và phản hồi cuối.

Bảng~\ref{tab:agent-attack-surface} tóm tắt các điểm tấn công chính trong
một agentic system.

\begin{table}[H]
\centering
\caption{Bề mặt tấn công chính trong hệ thống LLM agent.}
\label{tab:agent-attack-surface}
\begin{tabularx}{\textwidth}{
    >{\bfseries}p{0.22\textwidth}
    p{0.31\textwidth}
    X
}
\toprule
Thành phần
& Kênh tấn công
& Rủi ro chính \\
\midrule

User prompt
& Chỉ dẫn trực tiếp từ người dùng
& Jailbreak, direct prompt injection hoặc yêu cầu vượt chính sách. \\

External content
& Website, email, tài liệu, kết quả retrieval
& Indirect prompt injection, dữ liệu không đáng tin cậy bị hiểu nhầm là
instruction. \\

Memory
& Thông tin được lưu qua nhiều phiên
& Memory poisoning, ghi nhớ sai lệch hoặc lưu chỉ dẫn không đáng tin cậy. \\

Planner
& Kế hoạch trung gian của agent
& Phân rã nhiệm vụ theo hướng sai, bỏ qua ràng buộc hoặc chọn hành động
không phù hợp. \\

Tools
& API, database, file system, browser, code runner
& Tool misuse, rò rỉ dữ liệu, thao tác ngoài quyền hạn hoặc gọi công cụ sai
ngữ cảnh. \\

Tool output
& Dữ liệu trả về từ công cụ
& Nội dung độc hại hoặc chỉ dẫn gián tiếp được đưa ngược vào context. \\

Environment
& Website, ứng dụng, robot hoặc hệ thống vật lý
& Hành động có tác động bên ngoài, khó đảo ngược hoặc ảnh hưởng đến người
dùng khác. \\

\bottomrule
\end{tabularx}
\end{table}

Điểm cốt lõi là agent thường đặt nhiều loại dữ liệu vào cùng một context
window. Nếu hệ thống không phân tách rõ instruction, user request và
external content, mô hình có thể bị thuyết phục làm theo những chỉ dẫn
không đáng tin cậy. Đây là lý do agent security cần được thiết kế ở cấp
kiến trúc, không chỉ dựa trên hành vi từ chối của LLM.

\subsection{Prompt injection trực tiếp và gián tiếp}
\label{subsec:direct-indirect-prompt-injection}

Prompt injection là dạng tấn công trong đó attacker đưa chỉ dẫn vào context
nhằm làm thay đổi luồng điều khiển của ứng dụng LLM. Có thể phân biệt hai
dạng chính: direct prompt injection và indirect prompt injection.

\textbf{Direct prompt injection} xuất hiện khi người dùng trực tiếp đưa chỉ
dẫn đối nghịch vào prompt. Trong trường hợp này, attacker và user thường là
cùng một thực thể. Mục tiêu là khiến mô hình bỏ qua hoặc làm sai lệch các
chỉ dẫn cấp cao hơn, chẳng hạn system instruction hoặc application
instruction.

\textbf{Indirect prompt injection} xuất hiện khi chỉ dẫn đối nghịch không
nằm trong prompt trực tiếp của người dùng, mà nằm trong dữ liệu bên ngoài
được agent đọc hoặc truy xuất. Ví dụ, nội dung từ một trang web, tài liệu,
email hoặc tool output có thể chứa câu lệnh làm agent thay đổi kế hoạch.
Người dùng có thể không biết rằng dữ liệu bên ngoài chứa chỉ dẫn đó.

Greshake et al. chỉ ra rằng các ứng dụng tích hợp LLM làm mờ ranh giới
giữa dữ liệu và chỉ dẫn: nội dung được retrieval hoặc xử lý như dữ liệu
vẫn có thể ảnh hưởng đến hành vi điều khiển của ứng dụng
\cite{greshake2023indirect}. Đối với agent, vấn đề này nghiêm trọng hơn
chatbot vì agent có thể hành động dựa trên dữ liệu đã bị thao túng.

Có thể biểu diễn indirect prompt injection bằng luồng sau:

\begin{equation}
    u
    \rightarrow
    R(q)
    \rightarrow
    c_{\mathrm{ext}}
    \rightarrow
    F_{\theta}(u, c_{\mathrm{ext}}, s)
    \rightarrow
    a,
    \label{eq:indirect-prompt-injection-flow}
\end{equation}

trong đó $u$ là yêu cầu hợp lệ của người dùng, $R(q)$ là bước retrieval
hoặc tool call, $c_{\mathrm{ext}}$ là nội dung bên ngoài không đáng tin cậy,
$s$ là system instruction và $a$ là hành động của agent. Nếu
$c_{\mathrm{ext}}$ chứa chỉ dẫn đối nghịch nhưng được mô hình xử lý như
instruction, agent có thể thực hiện hành động ngoài ý định ban đầu.

\subsection{Tool-use attacks}
\label{subsec:tool-use-attacks}

Tool-use attacks nhắm vào khả năng agent gọi công cụ. Trong chatbot, phản
hồi không an toàn thường là văn bản. Trong agentic system, rủi ro có thể
nằm ở việc agent chọn sai công cụ, truyền sai tham số, truy xuất dữ liệu
không nên truy xuất hoặc thực hiện thao tác không phù hợp.

Một lời gọi công cụ có thể được mô tả bởi

\begin{equation}
    a_t = (\mathrm{tool}_t, \mathrm{args}_t),
    \label{eq:tool-call}
\end{equation}

trong đó $\mathrm{tool}_t$ là công cụ được chọn và $\mathrm{args}_t$ là tập
tham số truyền vào công cụ. Tool-use attack có thể tác động vào cả hai
thành phần này: làm agent chọn công cụ không phù hợp hoặc làm sai lệch tham
số hành động.

Các rủi ro thường gặp ở mức khái niệm gồm:

\begin{itemize}[leftmargin=1.2cm]
    \item agent gọi công cụ khi chưa có đủ quyền hoặc chưa cần thiết;
    \item agent truyền dữ liệu nhạy cảm vào công cụ không đáng tin cậy;
    \item agent dùng kết quả từ công cụ như instruction thay vì data;
    \item agent bỏ qua bước xác nhận người dùng trước hành động quan trọng;
    \item agent thực hiện chuỗi hành động đúng từng bước cục bộ nhưng sai
    mục tiêu tổng thể.
\end{itemize}

Do đó, phòng thủ cho tool-use không thể chỉ dựa vào output moderation.
Hệ thống cần kiểm tra trước khi gọi công cụ, giới hạn quyền theo nguyên tắc
least privilege, phân loại công cụ theo mức rủi ro và yêu cầu xác nhận với
các hành động có tác động bên ngoài.

\subsection{Memory poisoning}
\label{subsec:memory-poisoning}

Memory giúp agent duy trì thông tin qua nhiều lượt tương tác hoặc nhiều
phiên làm việc. Tuy nhiên, memory cũng tạo ra bề mặt tấn công mới. Nếu
attacker có thể làm agent ghi một thông tin sai lệch hoặc một chỉ dẫn không
đáng tin cậy vào bộ nhớ, thông tin đó có thể ảnh hưởng đến các quyết định
trong tương lai.

Có thể mô tả bộ nhớ dài hạn của agent tại thời điểm $t$ là $M_t$. Sau một
tương tác mới, agent cập nhật bộ nhớ bằng hàm ghi nhớ:

\begin{equation}
    M_{t+1}
    =
    \operatorname{Write}(M_t, e_t),
    \label{eq:memory-update}
\end{equation}

trong đó $e_t$ là thông tin được trích xuất từ tương tác hiện tại. Memory
poisoning xảy ra khi $e_t$ chứa thông tin không đáng tin cậy nhưng vẫn được
ghi vào $M_{t+1}$ và sau đó được retrieval trong một nhiệm vụ khác.

Rủi ro của memory poisoning nằm ở tính bền vững. Một prompt injection thông
thường có thể chỉ ảnh hưởng đến một phiên làm việc. Ngược lại, memory
poisoning có thể tạo ảnh hưởng dài hạn nếu agent tiếp tục truy xuất và tin
vào memory đã bị nhiễm. Vấn đề này đặc biệt quan trọng với các assistant cá
nhân hóa, web agents hoặc enterprise agents có bộ nhớ lâu dài.

Các biện pháp giảm thiểu ở mức hệ thống gồm:

\begin{itemize}[leftmargin=1.2cm]
    \item phân biệt rõ memory về sự kiện, sở thích và instruction;
    \item không cho phép external content tự động trở thành chỉ dẫn dài hạn;
    \item gắn metadata về nguồn, thời điểm và mức độ tin cậy cho memory;
    \item yêu cầu xác nhận trước khi ghi các memory có ảnh hưởng lớn;
    \item định kỳ kiểm tra, làm sạch hoặc hết hạn các memory rủi ro;
    \item dùng trust-aware retrieval để ưu tiên memory đáng tin cậy.
\end{itemize}

Nói cách khác, memory không nên được xem là vùng lưu trữ mặc định đáng tin
cậy. Nó cần có cơ chế kiểm soát quyền ghi, quyền đọc và độ tin cậy giống
như các thành phần dữ liệu khác trong hệ thống.

\subsection{Jailbreaking LLM-controlled robots}
\label{subsec:llm-controlled-robots}

Khi LLM được dùng để điều khiển robot hoặc hệ thống vật lý, rủi ro
jailbreak vượt ra ngoài văn bản. Một phản hồi không an toàn trong chatbot
có thể gây hại ở mức thông tin. Trong robot, hành vi không an toàn có thể
trở thành hành động trong môi trường thực. Vì vậy, đánh giá an toàn cần
xem xét cả lời gọi hành động, trạng thái robot, ràng buộc vật lý và cơ chế
dừng khẩn cấp.

RoboPAIR là một công trình tiêu biểu mở rộng jailbreak từ chatbot sang
LLM-controlled robots \cite{robey2024robopair}. Thay vì chỉ đánh giá xem mô
hình có sinh văn bản không an toàn hay không, hướng nghiên cứu này xem xét
liệu một hệ thống robot do LLM điều khiển có thể bị dẫn đến hành động trái
với chính sách an toàn hay không.

Trong bối cảnh robot, một action có thể được biểu diễn bởi

\begin{equation}
    a_t =
    \operatorname{Act}(p_t, o_t, r_t),
    \label{eq:robot-action}
\end{equation}

trong đó $p_t$ là kế hoạch, $o_t$ là quan sát môi trường và $r_t$ là tập
ràng buộc an toàn. Một hệ thống an toàn cần bảo đảm rằng action cuối cùng
không chỉ hợp lệ về mặt ngôn ngữ mà còn hợp lệ về mặt vật lý, quyền hạn và
bối cảnh.

Do đó, robot hoặc embodied agents cần nhiều lớp bảo vệ hơn chatbot:

\begin{itemize}[leftmargin=1.2cm]
    \item giới hạn không gian hành động mà LLM được phép đề xuất;
    \item dùng controller xác minh hành động trước khi thực thi;
    \item yêu cầu xác nhận con người với hành động rủi ro cao;
    \item áp dụng sandbox hoặc simulation trước khi hành động thật;
    \item dùng safety monitor độc lập với LLM;
    \item có cơ chế dừng khẩn cấp và rollback khi khả thi.
\end{itemize}

Điểm quan trọng là LLM không nên là lớp quyết định an toàn cuối cùng cho
hệ thống vật lý. Các ràng buộc an toàn cần được thực thi ở cấp controller
và môi trường, nơi có thể kiểm soát trực tiếp hành động.

\subsection{Agent-specific defenses}
\label{subsec:agent-specific-defenses}

Defense cho chatbot thường tập trung vào input guard, model alignment và
output moderation. Với agentic systems, các defense này vẫn cần thiết nhưng
chưa đủ. Agent cần các cơ chế bảo vệ bổ sung ở cấp context, memory, tool,
trajectory và permission.

Bảng~\ref{tab:agent-specific-defenses} tóm tắt các nhóm phòng thủ đặc thù
cho agent.

\begin{table}[H]
\centering
\caption{Các nhóm phòng thủ đặc thù cho agentic systems.}
\label{tab:agent-specific-defenses}
\begin{tabularx}{\textwidth}{
    >{\bfseries}p{0.25\textwidth}
    p{0.31\textwidth}
    X
}
\toprule
Nhóm defense
& Cơ chế chính
& Mục tiêu \\
\midrule

Context isolation
& Tách system instruction, user instruction và external content
& Giảm khả năng dữ liệu ngoài bị hiểu nhầm là instruction. \\

Tool permission
& Gán quyền truy cập công cụ theo nhiệm vụ, vai trò và mức rủi ro
& Ngăn agent gọi công cụ ngoài phạm vi cần thiết. \\

Action validation
& Kiểm tra công cụ và tham số trước khi thực thi
& Chặn hành động có tác động không mong muốn. \\

Human confirmation
& Yêu cầu xác nhận với hành động nhạy cảm hoặc không thể đảo ngược
& Giảm rủi ro khi agent không chắc chắn hoặc hành động có hậu quả lớn. \\

Memory access control
& Kiểm soát quyền ghi, đọc và xóa memory
& Giảm memory poisoning và retrieval từ nguồn không đáng tin cậy. \\

Trajectory monitoring
& Giám sát kế hoạch, tool calls và trạng thái trung gian
& Phát hiện rủi ro trước khi final response hoặc hành động cuối xảy ra. \\

Sandboxing
& Thực thi hành động trong môi trường giới hạn
& Giảm tác động khi agent hoặc tool bị thao túng. \\

Least privilege
& Chỉ cấp quyền tối thiểu cần thiết cho từng nhiệm vụ
& Giới hạn thiệt hại nếu một bước của agent bị compromise. \\

\bottomrule
\end{tabularx}
\end{table}

Một nguyên tắc quan trọng là agent không nên tự quyết định toàn bộ quyền
hạn của mình. Nếu LLM vừa lập kế hoạch, vừa chọn công cụ, vừa tự phê duyệt
hành động và vừa đánh giá an toàn, hệ thống dễ rơi vào single point of
failure. Thay vào đó, cần tách vai trò giữa LLM controller, permission
system, action validator và monitoring layer.

\subsection{Benchmark cho agent security}
\label{subsec:agent-security-benchmarks}

Benchmark cho chatbot thường đánh giá prompt và response. Benchmark cho
agent cần đánh giá task completion, tool-use correctness, security
violation, indirect prompt injection và khả năng duy trì utility dưới tấn
công. Do đó, benchmark agent security phức tạp hơn benchmark text-only.

AgentDojo là một benchmark tiêu biểu cho prompt injection attacks và
defenses trong LLM agents \cite{debenedetti2024agentdojo}. Benchmark này
đánh giá agent trong môi trường có công cụ và dữ liệu không đáng tin cậy,
nhấn mạnh cả utility và security. Điểm quan trọng là AgentDojo không chỉ là
một tập prompt tĩnh mà là môi trường động để thiết kế task, attack và
defense.

Agent Security Bench (ASB) mở rộng hướng đánh giá agent security bằng cách
formalize nhiều loại attack và defense trong các kịch bản agent khác nhau
\cite{zhang2024asb}. Các benchmark dạng này phản ánh nhu cầu đánh giá agent
ở nhiều bước, nhiều công cụ và nhiều loại hành vi, thay vì chỉ kiểm tra
response cuối.

Một benchmark agent security nên bao gồm ít nhất các thành phần sau:

\begin{enumerate}[leftmargin=1.2cm]
    \item \textbf{Benign tasks:}
    nhiệm vụ hợp lệ để đo task utility;

    \item \textbf{Adversarial tasks:}
    nhiệm vụ hoặc môi trường có injection để đo robustness;

    \item \textbf{Tool environment:}
    tập công cụ, quyền hạn và dữ liệu bên ngoài;

    \item \textbf{Security policy:}
    định nghĩa rõ hành vi nào được phép hoặc không được phép;

    \item \textbf{Trajectory logs:}
    ghi lại kế hoạch, tool calls, tool outputs và final response;

    \item \textbf{Scoring functions:}
    đo task success, security violation, over-defense, cost và latency.
\end{enumerate}

Khác với benchmark text-only, benchmark agent cần đánh giá cả hai mục tiêu:
agent phải hoàn thành nhiệm vụ hợp lệ và không bị hijack bởi dữ liệu không
đáng tin cậy. Nếu defense chặn mọi tool call để đạt security cao, agent có
thể mất utility. Nếu agent ưu tiên hoàn thành nhiệm vụ mà bỏ qua kiểm soát
an toàn, security violation có thể tăng.

\subsection{Tổng kết chương}
\label{subsec:agentic-systems-summary}

Chương này đã mở rộng bài toán jailbreak từ chatbot sang agentic systems.
Trong chatbot, rủi ro thường tập trung vào phản hồi văn bản. Trong agent,
rủi ro xuất hiện trên toàn bộ trajectory, bao gồm prompt, external content,
memory, plan, tool call, tool output, action và environment.

Các dạng rủi ro chính gồm direct prompt injection, indirect prompt
injection, tool-use attacks, memory poisoning và jailbreak đối với hệ thống
robot được điều khiển bởi LLM. Điểm chung của các rủi ro này là chúng khai
thác việc agent đưa dữ liệu không đáng tin cậy vào quá trình ra quyết định.

Do đó, phòng thủ cho agent cần vượt ra ngoài input/output moderation. Một
hệ thống agent an toàn cần context isolation, permission control, action
validation, memory access control, trajectory monitoring, sandboxing và
human confirmation cho hành động rủi ro cao. Đồng thời, evaluation cần đo
cả task utility và security violation trên toàn bộ trajectory.

Chương tiếp theo phân tích các công trình tiêu biểu được tutorial đề cập,
bao gồm các attack, defense và benchmark đại diện cho từng hướng nghiên
cứu.